\scanpage{104}

CORNU\qquad Philippe\qquad 16 Juin 1978
\mspage{I}

$C_4$\quad Mathématiques

\figure{\begin{tikzpicture}[x=1cm,y=1cm]\draw (0,0)--(6.8,.06);\end{tikzpicture}}

Ces deux projections stéréographiques de l'icosaèdre (par rapport a une face et par rapport à une arête) font apparaître quelques propriétés \unsure{Combinatoire} de ce polyèdre

Par construction, on est en présence de 15 cercles distincts dont 2 ont leur centre à l'infini pour la projection par rapport a une arête, et 3 pour l'autre projection

La première constatation que l'on peut faire est que chaque cercle porte deux arêtes distinctes. (on a donc bien $2\times15=30$ arêtes.

(De meme chaque cercle porte 4 sommets, correspondant aux deux arêtes, tous distincts, chaque sommet etant le lieu d'intersection de 5 cercles on a $\dfrac{4\times15}{5}=12$ sommets. Enfin, sur chaque cercle comme on l'a vu, se trouvent 2 arêtes, de part et d'autre desquelles se trouvent une face. Chaque face etant \unsure{limité} par 3 arêtes non \unsure{cocyclique} on a
\[
\frac{4\times15}{3}=20\text{ faces.})
\]

En meme temps, l'on peut voir que l'on est aussi en présence du dodécaèdre, polyèdre dual de l'icosaèdre dont le dual est l'icosaèdre.

\margin{103}
